\documentclass[10pt,a4paper]{amsart}
\usepackage[margin=19mm]{geometry}
\usepackage{amsmath,amssymb,mathtools}
\usepackage[scr=rsfs,cal=euler]{mathalfa}
\usepackage{xfrac}
\usepackage[unicode,hidelinks,bookmarks=false]{hyperref}
\newcommand{\ie}{i.e.\ }
\newcommand{\cf}{cf.\ }
\newcommand{\resp}{respectively\ }
\newcommand{\Subsection}{\subsection}
\providecommand{\nobreakdash}{\nobreak\hbox{-}}
\setlength{\emergencystretch}{2em}
\multlinegap0pt
\usepackage{amssymb}
\usepackage{mathtools}
\usepackage[shortlabels]{enumitem}
\newtheorem{lemma}{Lemma}
\newtheorem{proposition}{Proposition}
\newcommand{\Id}{I}
\newcommand{\R}{\mathbb R}
\newcommand{\T}{\mathbb T}
\newcommand{\calN}{\mathcal N}
\newcommand{\divv}{\operatorname{div}}
\newcommand{\eps}{\varepsilon}
\newcommand{\norm}[1]{\left\lVert #1\right\rVert}
\title{Residual-Work Compatibility Criteria and Defect-Measure Compactness\\ for Positive-Cone Viscoelastic Reynolds States}
\author{Sai Peng}
\date{Source: arXiv:2606.25309v2 (CC BY 4.0)}
\begin{document}
\maketitle

\noindent\emph{Source and licence.} Residual-Work Compatibility Criteria and Defect-Measure Compactness\\ for Positive-Cone Viscoelastic Reynolds States. Version arXiv:2606.25309v2 (CC BY 4.0), distributed by arXiv under the Creative Commons Attribution 4.0 International licence, http://creativecommons.org/licenses/by/4.0/, as recorded in the arXiv licence metadata for this exact version and shown on the abstract page https://arxiv.org/abs/2606.25309v2. Full licence notice: \texttt{licenses/arxiv-2606.25309-CC-BY-4.0.txt}.\par\smallskip

\noindent\textit{Editorial scope.} Counted unit: the article's proposition prop:concrete-layer (source0.tex lines 2144--2191). The article's complete original proof of it is delivered with it (source0.tex lines 2193--2320, one \texttt{proof} environment of 128 lines). No mathematics is written, repaired or re-derived: the statement and the proof are the article's own lines, in the article's own notation, and no retained line is substituted or re-flowed.  Retained uncounted as the source-local context the proof needs: lines 2110--2134.  Nothing was omitted from any retained span, so the package carries no omission record.  No retained line contains a citation, so the wrapper carries no bibliography and no bibliography entry.  No retained line contains a figure environment, an image-inclusion command or drawing code, so no image is delivered and the package declares no asset dependency.  Editorial formatting only, in addition: an AMS \texttt{amsart} preamble that restates the article's own declarations and macros used by the retained lines, namely the environments \texttt{lemma}, \texttt{proposition} on separate counters, together with the standard packages those lines need.  The retained environments print in this document as Lemma 1, Proposition 1 in order of appearance, whereas the article prints the same items with the numbering of its own full document.  The complete native source of the article as distributed in the arXiv e-print for this exact version is bundled unchanged as \texttt{sources/203592/source0.tex}.
\begin{lemma}[Finite-thickness scaling]
\label{lem:finite-thickness-scaling}
Let \(f\in C_c^\infty((-1/4,1/4))\), \(a\in\R\), and
\[
  f_\eps(x_2)=\eps^{-a}f(x_2/\eps)
\]
on \(\T^2\), for \(0<\eps\ll1\), by viewing the expression as a periodic
function of the coordinate \(x_2\).  Then
\[
  \norm{f_\eps}_{L^2(\T^2)}
  =
  \eps^{1/2-a}\norm{f}_{L^2(\R)},
  \qquad
  \norm{\partial_2 f_\eps}_{L^2(\T^2)}
  =
  \eps^{-1/2-a}\norm{f'}_{L^2(\R)}.
  \label{eq:finite-thickness-scale}
\]
If \(F\) is a smooth bounded function with \(F(0)=0\), then
\[
  \norm{F(f(x_2/\eps))}_{L^2(\T^2)}
  \le C_F\eps^{1/2}.
  \label{eq:bounded-layer-scale}
\]
\end{lemma}

\begin{proposition}[Concrete finite-thickness shear-layer obstruction]
\label{prop:concrete-layer}
Assume \(\nu>0\).  Fix a time window \(E\subset[0,T]\) of positive measure and
\(\beta>0\).  Let \(\theta,\chi\in C_c^\infty((-1/4,1/4))\), with
\[
  \int_{\R}\theta(y)\,dy=0,\qquad \theta'\not\equiv0,\qquad \chi\not\equiv0.
\]
Define
\[
  \theta_\varepsilon(x_2)=\theta(x_2/\varepsilon),\qquad
  \chi_\varepsilon(x_2)=\chi(x_2/\varepsilon),
\]
and set
\[
  H=\begin{pmatrix}1&0\\0&-1\end{pmatrix},\qquad
  u^\varepsilon(x)=
  \left(\varepsilon^{-\beta}\theta_\varepsilon(x_2),0\right),
\]
\[
  B^\varepsilon(x)=b\chi_\varepsilon(x_2)H,\qquad
  A^\varepsilon=e^{B^\varepsilon}.
\]
Let
\[
  R^\varepsilon=-2\nu D(u^\varepsilon)-\alpha(A^\varepsilon-\Id),
  \qquad
  S^\varepsilon=\mathcal S[u^\varepsilon,B^\varepsilon].
  \label{eq:concrete-defects}
\]
Then
\((u^\varepsilon,0,B^\varepsilon,R^\varepsilon,S^\varepsilon)\) is a smooth
positive-cone Reynolds state with fixed mean velocity.  Its residual-work data
obey
\[
  W_E^\varepsilon\ge c_W\varepsilon^{-(2\beta+1)},\qquad
  L_E^\varepsilon\le C_L\varepsilon^{1/2},
\]
\[
  S_E^\varepsilon\le C_S\varepsilon^{-(\beta+1/2)},\qquad
  \eta_E^\varepsilon\le1 .
\]
Thus the bulk exponents may be taken as
\[
  w=2\beta+1,\qquad \ell=0,\qquad r=\beta+\frac12,\qquad q=0,
\]
and the family is not energy--entropy admissible for all sufficiently small
\(\varepsilon\).
\end{proposition}

\begin{proof}
The support and zero-mean assumptions on \(\theta\) make
\(u^\varepsilon\) a smooth periodic divergence-free velocity with zero mean.
Since \(B^\varepsilon\) is symmetric, \(A^\varepsilon=e^{B^\varepsilon}\) is
positive definite.  More explicitly,
\[
  A^\varepsilon
  =
  \begin{pmatrix}
  e^{b\chi(x_2/\eps)}&0\\
  0&e^{-b\chi(x_2/\eps)}
  \end{pmatrix},
\]
so \(A^\varepsilon\), \((A^\varepsilon)^{-1}\), and all smooth functions of
them are uniformly bounded in \(L^\infty\), independently of \(\eps\).
Writing \(u^\varepsilon=(U^\varepsilon(x_2),0)\), one has
\[
  \divv(u^\varepsilon\otimes u^\varepsilon)=0,\qquad
  \divv(2D(u^\varepsilon))=\Delta u^\varepsilon.
\]
Therefore
\[
  \divv R^\varepsilon
  =
  -\nu\Delta u^\varepsilon-\alpha\divv(A^\varepsilon-\Id)
  =
  \calN[u^\varepsilon,B^\varepsilon],
\]
and the conformation residual is matched by the definition of
\(S^\varepsilon\).  Moreover \(D(u^\varepsilon)\) and
\(A^\varepsilon-\Id\) are symmetric, hence \(R^\varepsilon\) is a symmetric
representative of the zero-mean momentum residual.

The deformation tensor is purely off-diagonal:
\[
  \nabla u^\varepsilon
  =
  \begin{pmatrix}
  0&\partial_2U^\varepsilon\\
  0&0
  \end{pmatrix},
  \qquad
  D(u^\varepsilon)
  =
  \frac12
  \begin{pmatrix}
  0&\partial_2U^\varepsilon\\
  \partial_2U^\varepsilon&0
  \end{pmatrix}.
\]
Since \(A^\varepsilon-\Id\) is diagonal, its Frobenius product with
\(D(u^\varepsilon)\) vanishes:
\[
  (A^\varepsilon-\Id):D(u^\varepsilon)=0.
\]
Thus
\[
  P^\varepsilon(t)
  =
  \int_{\T^2}-R^\varepsilon:D(u^\varepsilon)\,dx
  =
  2\nu\int_{\T^2}|D(u^\varepsilon)|^2\,dx .
\]
Since
\[
  \partial_2U^\varepsilon
  =
  \varepsilon^{-\beta-1}\theta'(x_2/\varepsilon),
\]
and \(|D(u^\eps)|^2=\frac12|\partial_2U^\eps|^2\), Lemma~\ref{lem:finite-thickness-scaling}
gives
\[
  P^\eps(t)
  =
  \nu\norm{\partial_2U^\eps}_{L^2(\T^2)}^2
  =
  \nu\eps^{-(2\beta+1)}\norm{\theta'}_{L^2(\R)}^2.
\]
Thus
\[
  W_E^\eps
  =
  |E|\nu\norm{\theta'}_{L^2(\R)}^2
  \eps^{-(2\beta+1)}.
  \label{eq:W-exact-scale}
\]
The lever is
\[
  G^\varepsilon
  =
  \Id-(A^\varepsilon)^{-1}
  =
  \begin{pmatrix}
  1-e^{-b\chi(x_2/\eps)}&0\\
  0&1-e^{b\chi(x_2/\eps)}
  \end{pmatrix}.
\]
It is a bounded smooth function of \(b\chi(x_2/\eps)\), vanishing outside a
layer of thickness \(O(\eps)\).  Lemma~\ref{lem:finite-thickness-scaling} gives
\[
  L_E^\varepsilon\le C_L\eps^{1/2}.
  \label{eq:L-layer-scale}
\]
Finally,
\[
  \mathcal S[u^\varepsilon,B^\varepsilon]
  =
  -\nabla u^\varepsilon A^\varepsilon
  -A^\varepsilon(\nabla u^\varepsilon)^T
  +\lambda^{-1}(A^\varepsilon-\Id),
\]
because the state is steady and \(u^\varepsilon\cdot\nabla A^\varepsilon=0\).
The matrices \(-\nabla u^\varepsilon A^\varepsilon\) and
\(-A^\varepsilon(\nabla u^\varepsilon)^T\) are off-diagonal and contain the
factor \(\partial_2U^\varepsilon\) multiplied by the uniformly bounded
entries \(e^{\pm b\chi(x_2/\eps)}\).  Hence their \(L^2\) norm is bounded by
\[
  C\eps^{-(\beta+1/2)}.
\]
The relaxation term is a bounded smooth function of \(b\chi(x_2/\eps)\), so its
\(L^2\) norm is \(O(\eps^{1/2})\).  Therefore
\[
  S_E^\varepsilon\lesssim\varepsilon^{-(\beta+1/2)}.
\]
Using the weaker lever bound \(L_E^\varepsilon\le C_L\) corresponds to
\(\ell=0\).  Since \(2\beta+1>\beta+1/2\), the bulk power-law exclusion test
applies.
\end{proof}

\end{document}
